\clearpage
\SourceLeaf{377}%
\begin{center}\large [KONFORME ABBILDUNG. KRÜMMUNGSMASS.]\end{center}
\begin{center}[1.]\end{center}
\begin{center}\textsc{Gauss} an \textsc{Schumacher}. Göttingen, 5. Juli 1816.\end{center}

\mbox{. . . . . . .} Das Programm mit der Preisfrage Ihrer Gesellschaft habe ich noch nicht zu Gesicht bekommen. Mit \textsc{Lindenau} habe ich auch über eine Preisfrage gesprochen, die in der neuen Zeitschrift mit einem Preis von 100 Dukaten gestellt werden soll. Mir war eine interessante Aufgabe eingefallen, nämlich:
\begin{quote}\small
\noindent »\textbf{allgemein} eine gegebene Fläche so auf eine andere (gegebene) zu\\
»projizieren (abzubilden), dass das Bild dem Original in den kleinsten\\
»Teilen ähnlich wird«.
\end{quote}

Ein besonderer Fall liegt vor, wenn die erste Fläche eine Kugel und die zweite eine Ebene ist. Hier sind die stereografische Projektion und die Mercator-Projektion spezielle Lösungen. Gesucht ist aber die \textls[120]{allgemeine} Lösung, die alle speziellen Lösungen umfasst, und zwar für jede Art von Flächen.

\textsc{Monge} und \textsc{Ponsot} sollen darüber bereits im Journal philomathique gearbeitet haben (wie \textsc{Burckhardt} an \textsc{Lindenau} geschrieben hat). Da ich aber nicht genau weiß, wo, habe ich noch nicht nachsehen können und weiß daher nicht, ob die Lösungen jener Herren meiner Idee ganz entsprechen und die Sache erschöpfen. Andernfalls schiene mir dies einmal eine geeignete Preisfrage für eine Gesellschaft zu sein. Bei der hiesigen bin ich nur alle 12 Jahre an der Reihe, eine Preisfrage zu stellen. \mbox{. . . . . . .}
